\chapter{CLI e scaffolding}
\label{cap:19-cli-scaffolding}

Due strumenti stanno fra l'utente e la piattaforma: la CLI \code{nanofaas}, che
parla con un control plane, e \code{fn-init}, che genera lo scheletro di una
nuova funzione. Un terzo strumento --- quello che provvisiona l'infrastruttura
--- vive in un repository separato ed \`e trattato nel
\cref{cap:22-metodo-sperimentale}.

\section{Il confine fra gli strumenti}

La documentazione del progetto enuncia la divisione senza ambiguit\`a:

\begin{quote}\itshape
\code{nanofaas} \`e un client HTTP neutro rispetto al backend per un control
plane NanoFaaS. Non provvisiona infrastruttura, non chiama Helm o \code{kubectl},
e non ispeziona risorse specifiche del provider.
\end{quote}

La separazione ha una conseguenza di progetto: la CLI non ha alcuna dipendenza
da Kubernetes, da Docker o da qualunque backend, e la sua unica dipendenza
esterna \`e HTTP. \`E per questo che pu\`o essere compilata in un binario nativo
di dimensioni contenute.

\begin{center}
\small
\begin{adjustbox}{max width=\textwidth}
\begin{tabular}{@{}lp{8.2cm}@{}}
\toprule
\textbf{Obiettivo} & \textbf{Strumento} \\
\midrule
Costruire, registrare, invocare o ispezionare una funzione & CLI
\code{nanofaas} \\
Provvisionare una VM, installare k3s/Helm, distribuire immagini, eseguire E2E &
strumento di provisioning esterno \\
Eseguire il control plane in sviluppo & Gradle / Spring Boot \\
\bottomrule
\end{tabular}
\end{adjustbox}
\end{center}

\section{La CLI}

2\,426 righe di sorgente Java (\code{wc -l} su \code{clients/cli/src/main}, 2026-09-28), costruita su Picocli. La struttura dei comandi:

\begin{center}
\small
\begin{adjustbox}{max width=\textwidth}
\begin{tabular}{@{}ll@{}}
\toprule
\textbf{Comando} & \textbf{Ruolo} \\
\midrule
\code{fn apply} & Registra o aggiorna una funzione da un manifesto YAML. \\
\code{fn list} / \code{fn get} / \code{fn delete} & Ciclo di vita. \\
\code{fn test} & Invoca con un input di prova e confronta output e stato HTTP atteso. \\
\code{invoke} & Invocazione sincrona. \\
\code{enqueue} & Invocazione asincrona. \\
\code{exec get} & Stato di un'esecuzione. \\
\code{deploy} & Costruisce l'immagine e registra in un passo. \\
\bottomrule
\end{tabular}
\end{adjustbox}
\end{center}

La configurazione risolve in ordine: opzione da riga di comando, variabile
d'ambiente, file di configurazione (\code{\textasciitilde/.config/nanofaas/config.yaml}).
\`E la precedenza convenzionale, e l'assenza di sorprese \`e essa stessa una
propriet\`a desiderabile in uno strumento che compare dentro script di
esperimento.

\subsection{Due build, due scopi}

\begin{lstlisting}[language=bash]
./gradlew :nanofaas-cli:installDist      # distribuzione JVM
./gradlew :nanofaas-cli:nativeCompile    # eseguibile nativo GraalVM
\end{lstlisting}

La build JVM \`e quella di sviluppo. Quella nativa esiste perch\'e la CLI \`e
invocata dentro cicli: uno scenario di validazione la chiama decine di volte, e
il tempo di avvio della JVM --- dell'ordine del secondo --- si somma a ogni
chiamata. Un eseguibile nativo parte in decine di millisecondi.

\`E lo stesso argomento che il \cref{cap:20-build-packaging} applica al control
plane, ma con una differenza: per un processo di lunga durata il tempo di avvio
\`e ammortizzato e ci\`o che conta \`e il throughput a regime; per uno strumento a
riga di comando il tempo di avvio \emph{\`e} il tempo di esecuzione.

\subsection{Lo smoke test nativo}

\begin{lstlisting}[language=bash]
./gradlew :nanofaas-cli:nativeSmoke
\end{lstlisting}

Verifica help, versione e un comando HTTP contro uno stub locale. \`E una
verifica breve con uno scopo preciso: la compilazione nativa pu\`o riuscire e
produrre un binario che fallisce a runtime per una configurazione di riflessione
mancante --- l'errore caratteristico di GraalVM
(\cref{cap:20-build-packaging}). Un test che esercita almeno un percorso di
serializzazione reale \`e ci\`o che distingue ``ha compilato'' da ``funziona''.

Il rischio non \`e teorico. Il file di configurazione della riflessione della
CLI nominava undici tipi, e quattro comandi si rompevano nel binario nativo
mentre i loro test JVM passavano (\code{fn replicas get}, \code{control-plane
info}, \code{fn update -f} e \code{--{}-config <file>}: errori di lettura o di
serializzazione JSON e YAML). La correzione registra ogni \code{record} della CLI
pi\`u la classe di configurazione YAML (29 voci) e, soprattutto, aggiunge un test
(\code{NativeReflectConfigCoverageTest}) che \emph{scansiona} il codice
compilato alla ricerca dei \code{record} invece di elencarli: un tipo nuovo non
pu\`o pi\`u essere dimenticato. Con la vecchia configurazione il test nomina
esattamente i tipi mancanti. Dopo la correzione, secondo il messaggio del
commit, la CLI nativa e quella JVM producono lo stesso output per ogni comando
esercitato.

\subsection{Che cosa verifica \code{fn test}}

\code{fn test} passava un caso soltanto su uno stato 2xx, quindi un input che
descrive un percorso d'errore non poteva mai passare: un 422 deciso dalla
funzione arriva come HTTP 422 con stato di esecuzione \code{success} e il corpo
d'errore come output. Ogni scaffold di \code{fn-init} risponde 422 a un input
mancante, e dunque la CLI bocciava il \code{missing-input.json} dello scaffold
stesso appena generato. Ora un input pu\`o dichiarare \code{expectedStatusCode}
(chiave e predefinito 200 gi\`a usati dal corpus \code{correctness.json}): il caso
passa se lo stato HTTP coincide e l'output coincide con l'atteso, mentre un
errore di piattaforma continua a fallire. \`E la stessa distinzione fra errore
della funzione ed errore della piattaforma che il contratto della risposta
(\cref{cap:17-runtime-funzione}) fissa lato runtime.

\section{\code{fn-init}}

Genera lo scheletro di una nuova funzione: struttura di directory, file di
build, handler di esempio, Dockerfile, manifesto. \`E scritto in Python e vive
sotto \code{tools/}.

Il progetto ha registrato su questo strumento una regressione istruttiva: la
generazione dello scaffold JavaScript aveva perso la capacit\`a di incorporare
l'SDK, e la perdita non era stata rilevata perch\'e il test corrispondente era
rimasto orfano --- verificava una funzionalit\`a che non veniva pi\`u esercitata.
\`E il modo caratteristico in cui uno strumento di scaffolding degrada: il codice
generato continua a esistere e a sembrare corretto, e solo chi lo usa scopre che
manca qualcosa.

\begin{damisurare}
Non esiste una verifica che il codice generato da \code{fn-init} sia
effettivamente costruibile e registrabile per ogni linguaggio supportato. Un
cancello di questo tipo --- generare, costruire, registrare, invocare, per
ciascun template --- \`e realizzabile con gli strumenti gi\`a presenti.
\todomis{costruibilit\`a e invocabilit\`a degli scaffold generati, per linguaggio}
\end{damisurare}

\section{Il manifesto di funzione}

Il formato che \code{fn apply} consuma \`e la serializzazione YAML di
\code{FunctionSpec} (\cref{cap:05-contratti-dati}). Un esempio con le politiche
discusse nella Parte III:

\begin{lstlisting}[language=yamlnf, caption={Manifesto con controllo di
concorrenza a budget e offload reattivo.}]
name: word-stats
image: nanofaas/word-stats:1.0
executionMode: DEPLOYMENT
runtimeMode: HTTP
timeoutMs: 5000
concurrency: 8
queueSize: 100
resources:
  requests: { cpu: 0.5, memoryMiB: 256 }
  limits:   { cpu: 2.0, memoryMiB: 512 }
scalingConfig:
  strategy: INTERNAL
  minReplicas: 1
  maxReplicas: 6
  metrics:
    - { type: queue_depth, target: "5" }
  concurrencyControl:
    mode: BUDGETED
    targetLatencyMs: 10
    weight: 2.0
    minTargetInFlightPerPod: 1
    maxTargetInFlightPerPod: 8
offload:
  mode: pressure
  targetUrl: http://cloud-host:8080
\end{lstlisting}

Il manifesto \`e il punto in cui tutte le politiche descritte nella Parte III si
incontrano, ed \`e utile leggerlo come indice: \code{scalingConfig.strategy}
seleziona l'autoscaler (\cref{cap:11-autoscaler}),
\code{concurrencyControl.mode} il controllore di concorrenza
(\cref{cap:12-concurrency-control}), \code{offload} la politica di deviazione
(\cref{cap:13-offload}), e \code{queueSize} dimensiona la coda che l'ammissione
protegge (\cref{cap:10-sync-queue}).
